\documentclass[11pt]{article}
\usepackage[T1]{fontenc}
\usepackage{lmodern}
\usepackage{amsmath,amssymb,amsthm}
\usepackage[hmargin=1in,vmargin=.9in]{geometry}
\usepackage[hidelinks]{hyperref}
\hypersetup{
  pdftitle={Pair-Sum Periods in Three-Move Subtraction Games},
  pdfauthor={Haofan Zhang}
}

\newtheorem{theorem}{Theorem}
\newtheorem{lemma}[theorem]{Lemma}
\newtheorem{corollary}[theorem]{Corollary}
\newcommand{\Z}{\mathbb{Z}}
\newcommand{\N}{\mathbb{N}_0}

\title{Pair-Sum Periods in Three-Move Subtraction Games}
\author{Haofan Zhang}
\date{}

\begin{document}
\maketitle
\vspace{-.8em}

\begin{abstract}
For a normal-play subtraction game with moves $0<a<b<c$, we characterize
when each of $a+b$, $c+a$, and $c+b$ is a period from heap size zero of the
win--loss sequence. The end-of-block test of Mikl\'os and Post and the
initial-block formula of Manabe give a single translated-interval test for
$c+b$. We use reflection to relate this test to Manabe's admissible
angles and evaluate its endpoints to obtain the arithmetic test for
least period $c+b$. This gives self-contained proofs of Manabe's
Conjecture~57 and Mikl\'os--Post's Conjecture~3. We give the explicit
change of variables connecting the arithmetic test to Manabe's earlier
interval formulas and formalize both conjectures in Lean.
\end{abstract}

\section{Introduction}

The period of a subtraction game can depend on the move sizes and on
whether the pattern is required to repeat from heap size zero.
Flammenkamp \cite[Conjecture~6]{flammenkamp} asked whether the least
outcome period divides a pair sum in the non-additive three-move case.
Computational studies of three-move periods include the enumeration of
Benesh et al. \cite{benesh} and the conference report of Lacy et al.
\cite{lacy}. Zhang treated families including $b=2a$ and $b=a+1$
\cite[Propositions~3.1 and~4.1]{zhang}. Moriwaki treated the three
candidate periods for $b\le 2a$ with continuous subtraction parameters
\cite[Theorem~3]{moriwaki}.

Mikl\'os and Post gave a general test for repetition from zero, its
specialization to $c+b$, and a rule for removing moves longer than the
least pure outcome period. Their Conjecture~3 asks for an arithmetic
criterion for $c+b$ to be the least period
\cite[Proposition~7, Lemma~9, Corollary~9.1, and Conjecture~3]{miklos-post}.
Manabe gave sufficient tests on $c\bmod(a+b)$ and determined the
least periods in the covered cases. His Conjecture~57 asks for the
necessity of the tests at $c+a$ and $c+b$. His Remark~58 leaves the
explicit comparison with the arithmetic conjecture
\cite[Theorems~23 and~34, Conjecture~57, and Remark~58]{manabe}.

The $c+a$ necessity follows directly from Manabe's initial-block
calculation. For $c+b$, earlier criteria give the translated-interval
test, and Manabe's difference-set calculations give a route from that
test to his two avoidance conditions. We give a uniform reflection
proof of the equivalence (Section~\ref{sec:manabe}) and use the interval
endpoints and an elementary minimality argument to prove the arithmetic
criterion for all $0<a<b<c$ (Section~\ref{sec:arithmetic}). The latter
section also identifies the overlap with Manabe's earlier least-period
formulas. The companion Lean file proves the two conjectures in their
original scopes directly from the game recurrences.

\newpage
\section{Setup}

In a normal-play subtraction game a move removes an allowed number of
tokens from a heap, and a player with no legal move loses. A heap size is
a \emph{losing position} if the player to move loses under optimal play.
For any move set, the losing positions satisfy two rules. A losing
position has no legal move to another losing position, and every other
position has a legal move to a losing position. These rules uniquely
determine the losing set by induction on heap size. We call $p>0$ a
\emph{period} of the outcome if the losing indicator satisfies
$Y(n+p)=Y(n)$ for every $n\ge 0$.

Fix integers $0<a<b<c$. Write $W$ and $L$ for the losing sets of
$\{a,b\}$ and $\{a,b,c\}$, with indicators $z$ and $Y$, respectively.
Set $M=a+b$, $\delta=b-a$, and $\rho=c\bmod M$.
Intervals denote sets of integers. The losing residue set of the
two-move game is
\begin{equation}\label{eq:Z}
 Z=\{u\in[0,b):\lfloor u/a\rfloor\ \text{is even}\}
 \quad\text{in }\Z/M\Z .
\end{equation}
All expressions involving residues in $Z$ are read modulo $M$.
We use the gap sets
\[
 D=[a,b)\setminus Z,\qquad
 \Lambda=[\delta,b)\setminus Z .
\]
In particular $\Lambda\subseteq D$, since $[0,a)\subseteq Z$.
Formula~\eqref{eq:Z} is the usual two-move period block
\cite[Theorem~14]{miklos-post}. For completeness, the game below $b$
has only the move $a$, so the outcomes alternate in blocks of length
$a$. Every size in $[b,M)$ can subtract $b$ into $[0,a)$ and is
winning. The next lemma explains why this block repeats.

For $p$ larger than every move, the following formulation is equivalent
to the translating-zeros test of Mikl\'os and Post \cite[Lemma~9]{miklos-post}.
Manabe uses the same initial-block argument for three moves \cite[\S3]{manabe}.

\begin{lemma}[Repeating the initial block]\label{lem:repeat}
Let $p$ exceed every move in a subtraction game and let $H$ be the
losing positions in $[0,p)$. Then $H+p\N$ is the losing set if and only
if $H$ contains no two elements at distance $p-s$ for any allowed
move $s$.
\end{lemma}

\begin{proof}
If $x,y\in H$ and $y-x=p-s$, then $x+p$ and $y$ would both be losing
under repetition, with a legal move between them. Conversely, consider
the periodic set $H+p\N$. A move between members of this set would give
either distance $s$ within one block or distance $p-s$ between two
residues at the join. The first is excluded by the actual initial
outcomes and the second by hypothesis. If $n=r+kp$ is outside the
periodic set, with $0\le r<p$, some legal move at $r$ leads to $H$.
The same move at $n$ leads into $H+p\N$. The two rules for losing
positions now give the claim.
\end{proof}

For the two-move block of length $M$, the complementary distances
$M-a=b$ and $M-b=a$ are already excluded in the block. Hence
$z(n)$ is determined by $n\bmod M$ for every $n\ge 0$.
The two games agree below $c$. For $0\le u<b$ the first interval
affected by $c$ has the following formula from Manabe's initial-block calculation
\cite[\S3]{manabe}:
\begin{equation}\label{eq:tail}
 Y(c+u)=(1-z(u))(1-z(c+u-b)).
\end{equation}
Indeed, if $z(u)=1$, removing $c$ is a winning move. If $z(u)=0$,
then $u\ge a$ and $z(u-a)=1$ by the alternating block below $b$.
In this case the $a$-option $c+u-a$ is winning because it can remove
$c$ into $u-a$. The $c$-option $u$ is winning as well, and the
$b$-option is below $c$. This proves \eqref{eq:tail}. In particular
$[c,c+a)$ contains no losing positions.

\section{Exact tests for the three pair sums}

\begin{theorem}\label{thm:tests}
For all $0<a<b<c$, the following are necessary and sufficient
conditions for the indicated outcome periods:
\begin{align}
 a+b\quad &\Longleftrightarrow\quad
    Z\cap(\rho+Z)=\varnothing, \label{eq:short}\\
 c+a\quad &\Longleftrightarrow\quad
    (Z\cap[0,\delta))\cap\bigl(Z+(\delta-\rho)\bigr)=\varnothing,
    \label{eq:middle}\\
 c+b\quad &\Longleftrightarrow\quad
    \rho+\Lambda+a\subseteq Z. \label{eq:long}
\end{align}
Moreover, if the outcome has any period $p\le\max\{c,M\}$, then
$L=W$. If $L\ne W$, any period less than $2c$ is the least period.
\end{theorem}

Mikl\'os and Post showed that removing the moves longer than the least
pure outcome period leaves the sequence unchanged
\cite[Proposition~7]{miklos-post}. This gives the redundancy of $c$ when
the least period is less than $c$. The argument below uses the initial
block directly. Equation~\eqref{eq:short} is
Manabe's inert-move test \cite[Definition~19 and Proposition~25]{manabe},
stated here for the outcome sequence.

\begin{proof}
Suppose first that $p\le\max\{c,M\}$ is a period of $Y$. If
$p\le c$, then $Y=z$ throughout $[0,p)$. If $c<p\le M$, the
same agreement holds since $[c,p)$ lies in $[c,c+a)$ and is
also disjoint from $W$ by \eqref{eq:Z} and $b<c$. For an arbitrary
$n=r+kp$, $0\le r<p$, the following two checks show that $L$ is
the losing set of the base game. First, there is no move by $a$ or $b$
from one member of $L$ to another, since these are moves in
$\{a,b,c\}$. Second, if $n\notin L$, then $r\notin W$, so some
base move at $r$ reaches a member of $W\cap[0,p)=L\cap[0,p)$.
By periodicity this move at $n$ reaches $L$. Thus $L=W$.

The set $W$ satisfies the rules for the three-move game exactly when
it contains no two positions at distance $c$. Every position outside
$W$ already has a winning base move. By $M$-periodicity the distance
condition is exactly $Z\cap(\rho+Z)=\varnothing$. This proves
\eqref{eq:short}. For a purely periodic sequence its least period
divides every period, by the Euclidean algorithm on periods.
If $L\ne W$, a period $p<2c$ with a smaller least period would
therefore have a period at most $p/2<c$, a contradiction.

For $p=c+a$, we use the initial block and collision calculation of
\cite[Theorem~23]{manabe}. The initial losing block is $W\cap[0,c)$ by
\eqref{eq:tail}. The distances $p-a=c$ and $p-c=a$ are already
excluded since they are legal moves. A pair at distance $p-b=c-\delta$
has the form $x,x+c-\delta$ for $0\le x<\delta$. Both positions
are losing exactly when $x\in Z$ and $x+\rho-\delta\in Z$.
Lemma~\ref{lem:repeat} gives \eqref{eq:middle}.

For $p=c+b$, the distances $p-b=c$ and $p-c=b$ are legal moves.
A pair at the remaining distance $p-a=c+\delta$ must start at
$x\in[0,a)$, and each such $x$ is losing. Thus the test in
Lemma~\ref{lem:repeat} reduces to $Y(c+u)=0$ for all
$u\in[\delta,b)$. This is the end-of-block criterion of
Mikl\'os and Post \cite[Corollary~9.1]{miklos-post}.
By \eqref{eq:tail}, the requirement is automatic for $u\in Z$,
and otherwise it is $z(c+u-b)=1$. The latter residue is
$\rho+u+a$, which proves \eqref{eq:long}.
\end{proof}

\section{Manabe's admissible-angle conjecture}\label{sec:manabe}

Manabe's condition T1 is \eqref{eq:middle}. His condition T2 consists
of two avoidance conditions \cite[Definition~19 and Lemma~22]{manabe}.
\begin{equation}\label{eq:T2}
 (\rho+\Lambda)\cap Z=\varnothing,\qquad
 (\rho+D)\cap\Lambda=\varnothing
\end{equation}
We prove their equivalence with \eqref{eq:long} by reflection.

First, in $\Z/M\Z$ we have
\begin{equation}\label{eq:lambda}
 \Lambda=(\Z/M\Z)\setminus\bigl(Z\cup(Z-a)\bigr).
\end{equation}
For $u\in[0,\delta)$, if $u$ is outside $Z$, adding $a$ flips its
block parity while remaining below $b$, and hence enters $Z$.
For $u\in[b,M)$, adding $a$ wraps into $[0,a)$.
For $u\in[\delta,b)$, adding $a$ instead lands in $[b,M)$.
This recovers the identity in \cite[Lemma~22]{manabe}. Also
$Z\cap(Z-a)=\varnothing$, since two base losing positions cannot differ by $a$.

We now exhibit a reflection preserving $Z$. Write $b=qa+r$ with
$0\le r<a$. If $q$ is odd, $Z$ consists of equal blocks with
endpoints ranging from $0$ to $qa-1$, and $u\mapsto qa-1-u$
reverses the blocks. If $q$ is even, $u\mapsto (q-1)a-1-u$
reverses the full blocks and preserves the last block $[qa,b)$
modulo $M$. In the second case, on that last block the map is
$u\mapsto qa+b-1-u$ after adding $M$. Let the corresponding map in
either case be $u\mapsto R-u$.
Then the reflection
\[
 \iota(u)=R-a-u
\]
interchanges $Z$ and $Z-a$, so it preserves $\Lambda$ by
\eqref{eq:lambda}.

Now \eqref{eq:long} gives $\rho+\Lambda\subseteq Z-a$, which
implies the first condition of \eqref{eq:T2}. Applying $\iota$
gives $\Lambda-\rho\subseteq Z$. Since $D$ is disjoint from
$Z$, this is enough for the second condition of \eqref{eq:T2}.
Conversely, \eqref{eq:T2} and $\Lambda\subseteq D$ show that
$\rho+\Lambda$ avoids both $Z$ and $\Lambda$, so by
\eqref{eq:lambda} it lies in $Z-a$. This is \eqref{eq:long}.

The earlier difference-set formulas give another route. Using
\eqref{eq:lambda}, the forbidden residues in \eqref{eq:long} form
$(Z-\Lambda)\cup(\Lambda-\Lambda)$, and T2 forbids
$(Z-\Lambda)\cup(\Lambda-D)$. For $a\nmid b$, the arcs listed in
\cite[Proposition~29 and Appendix~A]{manabe} give the same union.
In his notation $b-a=\eta a+r$. When $\eta$ is even, the indexed arcs
of $\Lambda-D$ lie in $Z-\Lambda$ and the remaining arc is
$\Lambda-\Lambda$. When $\eta$ is odd, the extra arc of $\Lambda-D$
lies in $(\Lambda-\Lambda)\cup[1,a-1]$, and $[1,a-1]$ lies in
$Z-\Lambda$. The inclusion $\Lambda\subseteq D$ gives the reverse
inclusion of the unions.
When $b/a$ is an odd integer, $\Lambda$ is empty. When $b/a$ is an
even integer, fitting the translated length-$a$ interval in a
length-$a$ block forces $a\mid\rho$, which the primitive hypotheses
exclude. Manabe excludes the same case for T2
\cite[Theorem~30 and Proposition~61]{manabe}.

\begin{corollary}
Conjecture~57 of Manabe \cite[\S6]{manabe} holds. If the Sprague--Grundy
sequence is periodic from zero with least period $p=c+a$ or $p=c+b$,
the residue $\rho$ satisfies the test for that period.
The same necessity holds for the candidate $p=a+b$.
\end{corollary}

\begin{proof}
Every period of the Sprague--Grundy sequence is a period of the
outcome, since its value is zero exactly at losing positions.
Apply Theorem~\ref{thm:tests} and the equivalence just proved.
\end{proof}

These arguments apply to the specified periods. Restricting the
set of possible least periods remains a separate question
\cite[Conjecture~56]{manabe}.

\section{The arithmetic conjecture of Mikl\'os and Post}\label{sec:arithmetic}

Define
\[
 \rho=2aj+t,\qquad \sigma=(c-a)\bmod M=2aj'+t',
 \qquad 0\le t,t'<2a,
\]
where the residues $\rho,\sigma$ lie in $[0,M)$. As above,
$b=qa+r$ with $0\le r<a$.

For a purely periodic sequence, any eventual period is a period
from zero. Indeed, translate the two positions being compared by a
sufficiently large multiple of a pure period. Thus the least period
from zero agrees with the least eventual period when the preperiod is zero.

\begin{theorem}\label{thm:arithmetic}
The outcome has least period $c+b$ from zero if and only if the
applicable condition below holds.
\begin{enumerate}
 \setlength{\itemsep}{0pt}
 \item If $q$ is even, then
 $\sigma>0$, $t'\le r$, $2j'\le q$, and
 $t'+a\le 2r$ when $2j'=q$.
 \item If $q$ is odd and $r>0$, then
 $r\le t\le a$, with $t<a$ when $j=0$.
 \item If $q$ is odd and $r=0$, then $t\ne a$.
\end{enumerate}
For $a>1$ this is Conjecture~3 of Mikl\'os and Post \cite[\S7]{miklos-post}.
The last even-case inequality is their $t'\le 2r-a$ over the integers.
\end{theorem}

\begin{proof}
Since $M<c+b<2c$, the period $c+b$ is least exactly when it
exists and $L\ne W$, by Theorem~\ref{thm:tests}.
In terms of the base residues the two requirements are
\[
 \rho+\Lambda+a\subseteq Z,\qquad Z\cap(\rho+Z)\ne\varnothing .
\]
We evaluate these using the parity of $q$.

\smallskip
\noindent\textit{$q$ even.}
Write $q=2h$. The set $Z$ has full blocks
$[2ia,(2i+1)a)$ for $0\le i<h$ and a last block $[qa,b)$
of length $r$. The gap $\Lambda$ is $[(q-1)a+r,qa)$.
Thus $\rho+\Lambda+a$ is the cyclic interval starting at
$\sigma$ with length $a-r$. It must fit into a single block
of $Z$, since the blocks are separated by nonempty gaps,
including at the wrap. This gives $t'\le r$ in a full block
and $t'\le 2r-a$ in the last block, with $j'\le h$.

Among these possibilities, $\sigma=0$ is excluded because
$\rho=a$ translates $Z$ disjointly from itself. For
$\sigma>0$, the fit bounds give
$\rho=(2j'+1)a+t'<M$. The integer
\[
 z_*=(q-2j')a+(r-t')
\]
is in $Z$. If $j'>0$, it is in a full block because
$0\le r-t'<a$. If $j'=0$, then $t'>0$ and
$0\le r-t'<r$, so it is in the last block. Since
$\rho+z_*=M$ and $0\in Z$, the required overlap follows.

\smallskip
\noindent\textit{$q$ odd and $r>0$.}
Write $q=2h+1$. The blocks of $Z$ are $[2ia,(2i+1)a)$ for
$0\le i\le h$, and $\Lambda=[qa,b)$.
The interval $\rho+\Lambda+a$ is $[\rho-r,\rho)$ modulo $M$.
Fitting into one block gives $2ia+r\le\rho\le(2i+1)a$
for some $0\le i\le h$. This is equivalent to $r\le t\le a$,
since $\rho<M$ excludes $j>h$ when $t\ge r$.
When $t<a$, $\rho$ itself lies in $Z$, giving overlap with
$\rho+0$. When $t=a$ and $j>0$, the residue of
$\rho+(q-1)a$ is $(2j-1)a-r$, which belongs to the
block $[2(j-1)a,(2j-1)a)$. As $(q-1)a\in Z$, this
gives overlap again. The excluded endpoint $j=0$, $t=a$
has $\rho=a$.

\smallskip
\noindent\textit{$q$ odd and $r=0$.}
Here $\Lambda$ is empty and $M$ is a multiple of $2a$.
The base set is the lift of $[0,a)$ modulo $2a$, so disjointness
from its translate occurs exactly when $t=a$. This proves the
last condition.
\end{proof}

\paragraph{Comparison with the earlier residue intervals.}
For primitive non-additive triples with $a\nmid b$, Manabe's
Theorem~30 gives a set $C_{a,b}$ of angles, and Theorem~34,
Proposition~48, and Corollary~49 give least outcome period $b+c$
throughout this set \cite{manabe}. In his notation $\eta=q-1$ and
$\varepsilon=r$. For $q=2h+1$, the interval $[r,a)$ is the $j=0$
case of our test. The other intervals are $[2ja+r,(2j+1)a]$ for
$1\le j\le h$. For $q=2h$, translating the allowed values of
$\sigma$ by $a$ gives $[a+1,a+r]$, then
$[(2j'+1)a,(2j'+1)a+r]$ for $1\le j'<h$, and finally
$[(q+1)a,qa+2r]$ when nonempty. These are exactly the intervals of
$C_{a,b}$, so the earlier least-period result gives the sufficient
direction in this range by the change of variables. For the reverse
direction, a least outcome period of $b+c$ forces T2 by the interval
criterion and the equivalence in Section~\ref{sec:manabe}. It rules
out the inert angles, for which $a+b$ is already a period.

The ingredients for the remaining ranges are available separately.
For primitive triples with $a>1$ and $a\mid b$, the admissible-angle
results and the necessity above give least period $b+c$ in the odd
quotient case and exclude it in the even case
\cite[Theorems~30 and~34, Corollary~49, and Proposition~61]{manabe}. This agrees
with the arithmetic test. In the odd case $t=a$ would imply $a\mid c$,
and in the even case the test requires $t'=0$, also forcing $a\mid c$.
In the primitive additive case with $a>1$,
$\rho=0$, $\sigma=b$, and $r>0$, so the arithmetic test fails. The
additive formula gives $a+2b-r$ for odd $q$ and $a(2b+r)$ for even
$q$, both different from $b+c$ \cite[Theorem~36]{manabe}.
Mikl\'os and Post's dilation rule reduces nonprimitive triples to
primitive ones. Dividing the move sizes by their common divisor
divides the remainders by that divisor and leaves the quotients fixed.
Their $a=1$ classification gives least pure period $b+c$ precisely
when $b$ is odd and $c$ is even, or when $b$ is even and $\rho$ is odd
and greater than $1$. These are the arithmetic conditions with $a=1$
\cite[Proposition~5 and Theorem~20]{miklos-post}. The proof above handles
all these cases by the same interval and minimality criteria.

\section*{Verification}

\paragraph{Computational check.}
We compared the criteria with the direct outcome recurrence using
the companion script \texttt{verify\_subtraction\_periods.py}.
All comparisons agree for the $161{,}700$ triples with
$1\le a<b<c\le100$, including the minimality claim for $c+b$.

\paragraph{Formal verification.}
The accompanying file \texttt{SubtractionPeriods.lean} formalizes both
conjectures in Lean 4.19.0 using the standard library.
\begin{itemize}
 \setlength{\itemsep}{0pt}
 \item \texttt{manabe\_conjecture\_57\_canonical} derives both avoidance
 implications from the canonical mex recurrence.
 \item \texttt{miklos\_post\_conjecture\_3\_original} derives the
 arithmetic equivalence for zero preperiod and least eventual period
 $b+c$ from the unseeded losing-position recurrence.
\end{itemize}
Manabe's own Lean supplement
formalizes his sufficient criteria and least-period results and lists
Conjecture~57 among the statements left unformalized
\cite{manabe-lean}. For both theorems in the present file, Lean reports
the axioms \texttt{propext}, \texttt{Classical.choice}, and \texttt{Quot.sound}.

\paragraph{Use of generative AI.}
OpenAI Codex developed the proof strategy and derivations, wrote the
verification script and Lean formalization, and prepared this draft
for human mathematical review.

\begin{thebibliography}{9}
\raggedright
\bibitem{benesh}
Bret Benesh, Jamylle Carter, Deidra A. Coleman, Douglas G. Crabill,
Jack H. Good, Michael A. Smith, Jennifer Travis, and Mark Daniel Ward,
\emph{Periods in Subtraction Games},
29th International Conference on Probabilistic, Combinatorial and
Asymptotic Methods for the Analysis of Algorithms, LIPIcs \textbf{110}
(2018), 8:1--8:3.
\href{https://doi.org/10.4230/LIPIcs.AofA.2018.8}{doi:10.4230/LIPIcs.AofA.2018.8}.

\bibitem{flammenkamp}
Achim Flammenkamp,
\emph{Lange Perioden in Subtraktions-Spielen},
Ph.D. thesis, Universit\"at Bielefeld, 1997.
\href{https://wwwhomes.uni-bielefeld.de/achim/diss.ps.gz}{Author's copy}.

\bibitem{lacy}
Kali Lacy, Bret Benesh, Jamylle Carter, Deidra Coleman, Douglas Crabill,
Jack Good, Michael Smith, Jennifer Travis, and Mark Ward,
\emph{A Computational Approach to the Structure of Subtraction Games},
Joint Statistical Meetings, 2019,
\href{https://ww2.amstat.org/meetings/jsm/2019/onlineprogram/AbstractDetails.cfm?abstractid=305112}
{poster abstract 305112}.

\bibitem{manabe}
Hikaru Manabe,
\emph{Purely Periodic Three-move Subtraction Games},
preprint, \href{https://arxiv.org/abs/2609.05358v2}{arXiv:2609.05358v2},
2026.

\bibitem{manabe-lean}
Hikaru Manabe,
\emph{Lean formalization of ``Purely Periodic Three-move Subtraction Games''},
\href{https://arxiv.org/src/2609.05358v2/anc/lean-verification/README.md}
{ancillary README to arXiv:2609.05358v2}, 2026.

\bibitem{miklos-post}
Istv\'an Mikl\'os and Logan Post,
\emph{Superpolynomial period lengths of the winning positions in the
subtraction game},
International Journal of Game Theory \textbf{53} (2024), 1275--1313,
\href{https://doi.org/10.1007/s00182-024-00911-5}{doi:10.1007/s00182-024-00911-5}.

\bibitem{moriwaki}
Yuto Moriwaki,
\emph{Subtraction Nim with Continuous Parameters},
preprint, \href{https://arxiv.org/abs/2606.11658}{arXiv:2606.11658}, 2026.

\bibitem{zhang}
Shenxing Zhang,
\emph{On the linearity of the periods of subtraction games},
Theoretical Computer Science \textbf{985} (2024), 114350,
\href{https://doi.org/10.1016/j.tcs.2023.114350}{doi:10.1016/j.tcs.2023.114350}.
\end{thebibliography}

\end{document}
